\section{LangGraph 实现 ReAct Agent}

\subsection{基于图的架构}

LangGraph 的 \texttt{create\_react\_agent} 实现更为现代化。只需定义模型和工具即可：

\begin{importantbox}{理解 LangGraph 的关键：虚边（条件边）}
LangGraph 的 ReAct Agent 包含两个核心节点：
\begin{itemize}
    \item \textbf{Agent 节点}：负责产生工具调用请求
    \item \textbf{Tools 节点}：负责执行工具并返回结果
\end{itemize}
Agent 节点上有一条\textbf{条件边}，决策是否要停止调用工具直接输出答案，还是继续调工具。
\end{importantbox}

\subsection{基于消息列表的管理}

与 LangChain 不同，LangGraph 完全基于\textbf{多轮对话消息}的方式管理交互：

消息流：Human Message $\rightarrow$ AI Message（含 tool\_calls） $\rightarrow$ Tool Message(s) $\rightarrow$ AI Message（最终回答）

\begin{knowledgebox}{并行工具调用}
LangGraph 支持\textbf{并行工具调用}——AI Message 中的 \texttt{tool\_calls} 是一个列表，可以同时调用多个工具。例如查询"澳网冠军的家乡"时，模型可能并行调用两次搜索工具。
\end{knowledgebox}

\subsection{本章小结}
LangGraph 的 ReAct 实现基于图结构和多轮消息列表，架构更清晰，支持并行工具调用。

